% ============================================================
% Chapitre 2 : Spécification des besoins
% ============================================================
\chapter{Spécification des besoins}

\begin{spacing}{1.2}
\minitoc
\thispagestyle{MyStyle}
\end{spacing}
\newpage

\section*{Introduction}

Le chapitre précédent a fixé les objectifs du projet. Ce chapitre les traduit en spécifications exploitables. Nous identifions d'abord les acteurs qui interagissent avec la plateforme, car c'est de leurs attentes respectives que découlent les fonctionnalités attendues. Nous formulons ensuite les besoins fonctionnels, puis les besoins non fonctionnels qui fixent les critères de qualité auxquels le système doit satisfaire. Le chapitre se termine par les cas d'utilisation, présentés sous forme de diagrammes puis détaillés pour les trois plus déterminants d'entre eux.

\section{Identification des acteurs}

Un acteur désigne une entité extérieure au système qui interagit avec lui pour en obtenir un service. Nous distinguons les acteurs principaux, qui utilisent directement la plateforme, des acteurs secondaires, qui sont des systèmes tiers sollicités par elle.

\subsection{Acteurs principaux}

Deux acteurs principaux ont été retenus : le client et l'utilisateur back-office. Ce dernier regroupe les profils internes de l'opérateur ; ses responsabilités de supervision, d'administration et de prise en charge des escalades sont précisées par la matrice des droits présentée au chapitre suivant.

\begin{itemize}
  \item \textbf{Le client.} Abonné de l'opérateur, disposant d'un compte sur la plateforme. Il peut créer ce compte depuis le portail et vérifier son identité auprès des services de l'opérateur, ce qui lui permet ensuite d'agir avec l'agent. Il s'adresse à la plateforme pour obtenir une information, faire réaliser une opération sur sa ligne ou signaler un problème, et accède par ailleurs à un portail où il retrouve sa situation, ses demandes et l'historique de ses échanges.
  \item \textbf{L'utilisateur back-office.} Il rassemble les moyens internes de supervision et de gouvernance du système : observation de l'activité, inspection des décisions, gestion des disponibilités, administration des règles métier et de la base documentaire, traitement des demandes escaladées et contrôle de l'intégrité des enregistrements. La matrice des droits du chapitre suivant distingue les trois rôles réels qui le composent : l'administrateur, le superviseur et le conseiller.
\end{itemize}

\subsection{Acteurs secondaires}

La plateforme ne détient pas les données métier de l'opérateur : elle les consulte et agit sur elles au travers de systèmes tiers, qui constituent autant d'acteurs secondaires.

\begin{itemize}
  \item \textbf{Le système de gestion de tickets.} Il demeure la référence pour tout ce qui concerne les tickets. La plateforme y crée, consulte, met à jour et clôture les tickets des clients.
  \item \textbf{Les systèmes métier de l'opérateur.} Ils portent la facturation et les soldes, la supervision du réseau et le provisionnement des services. La plateforme les interroge pour construire le contexte du client et leur adresse les actions approuvées.
  \item \textbf{Le service de notification.} Il achemine les confirmations et les avis destinés aux clients comme aux conseillers.
\end{itemize}

Le tableau \ref{tab:acteurs} récapitule le rôle de chaque acteur et son mode d'interaction avec la plateforme.

\begin{table}[H]
\centering
\renewcommand{\arraystretch}{1.25}
\begin{tabular}{|p{2.6cm}|p{2.2cm}|p{4.4cm}|p{4.6cm}|}
\hline
\textbf{Acteur} & \textbf{Type} & \textbf{Responsabilité principale} & \textbf{Mode d'interaction} \\
\hline
Client & Principal & Exprimer une demande, obtenir sa résolution ou souscrire à un service & Conversation vocale ou textuelle, portail client \\
\hline
Utilisateur back-office & Principal & Superviser l'activité, administrer les règles et la base documentaire, traiter les demandes escaladées, contrôler l'intégrité & Console d'administration et de supervision \\
\hline
Système de gestion de tickets & Secondaire & Détenir la référence des tickets & Interface de programmation applicative \\
\hline
Systèmes métier de l'opérateur & Secondaire & Détenir les données de facturation, de réseau et de provisionnement & Interfaces de service internes \\
\hline
Service de notification & Secondaire & Acheminer les confirmations et les avis & Interface de service interne \\
\hline
\end{tabular}
\caption{Acteurs de la plateforme et modes d'interaction}
\label{tab:acteurs}
\end{table}

\section{Besoins fonctionnels}

Les besoins fonctionnels décrivent ce que le système doit faire. Nous les avons regroupés par domaine et identifiés individuellement, afin que chacun puisse être suivi jusqu'au chapitre de réalisation.

\subsection{Conduite de la conversation}

La plateforme doit permettre au client d'engager un échange parlé ou écrit et de le mener à son terme sans contrainte de formulation. Elle doit conserver le fil de la conversation d'un tour de parole au suivant, de sorte que le client n'ait jamais à répéter une information déjà donnée. Elle doit identifier l'intention exprimée et, lorsque la demande demeure ambiguë, conduire une clarification bornée. Elle pose une question précise, adapte sa stratégie (question fermée ou épellation de la référence) si la réponse reste inexploitable, et au delà de trois tentatives infructueuses considère que c'est la clarification qui a échoué, et non le client, puis propose le passage à un conseiller. Elle doit accepter d'être interrompue à tout moment et reprendre la main sans perdre l'état de l'échange. Elle doit enfin pouvoir changer de langue en cours de conversation à la demande du client.

\subsection{Construction du contexte client}

La plateforme doit reconstituer automatiquement la situation du client au moment où il s'adresse à elle : son profil, son abonnement, sa situation de facturation, ses demandes récentes et l'historique de ses échanges antérieurs. Cette reconstitution doit être conditionnelle plutôt que systématique, chaque élément n'étant sollicité que lorsque la nature de la demande le justifie, afin de ne pas alourdir inutilement le traitement.

\subsection{Ancrage documentaire des réponses}

La plateforme doit répondre aux questions documentaires en s'appuyant sur la documentation propre à l'opérateur, et non sur les connaissances générales du modèle de langue. Lorsque cette documentation ne contient pas de réponse à la question posée, le système doit le signaler explicitement au client plutôt que produire une réponse vraisemblable. La base documentaire doit pouvoir être alimentée et mise à jour par l'administrateur sans intervention sur le code.

\subsection{Traitement des actions sensibles}

Toute opération engageant le compte du client doit être précédée d'une vérification d'identité, puis soumise à un jeu de règles métier produisant un verdict explicite. Ce verdict doit prendre l'une des trois valeurs suivantes : autorisée, refusée, ou à transmettre à un conseiller. Il doit être accompagné de l'identifiant de la règle qui l'a motivé et d'une justification lisible. Une action autorisée doit être exécutée une fois et une seule, y compris en cas de reprise ou de tentative concurrente. Le résultat doit être confirmé verbalement au client.

\subsection{Gestion des tickets}

La plateforme doit pouvoir consulter les tickets existants d'un client, en créer de nouveaux, les mettre à jour et les clôturer lorsque le problème est résolu. Elle doit connaître l'historique des tickets du client afin d'éviter la création de doublons et de tenir compte des difficultés déjà signalées. Elle doit en outre permettre au client d'ouvrir un litige sur une facture : l'ouverture retire la facture du règlement pendant l'instruction, et sa clôture restaure le statut antérieur.

\subsection{Escalade et rappel}

La plateforme doit reconnaître les situations qui appellent une intervention humaine, qu'elles résultent d'une règle métier, d'une demande explicite du client ou d'un mécontentement persistant. Lorsqu'un conseiller est disponible, elle doit lui transférer la conversation. Dans le cas contraire, elle doit proposer un rappel, négocier un créneau compatible avec les disponibilités réelles des conseillers, réserver ce créneau et en informer le client. Dans les deux cas, le conseiller concerné doit recevoir le dossier de l'échange.

\subsection{Espace client}

Le client doit disposer d'un espace personnel lui permettant de consulter son profil, sa situation de facturation, ses demandes, ses rappels programmés et l'historique de ses conversations. Il doit également pouvoir consulter et effacer les informations que la plateforme conserve de ses échanges antérieurs, et n'accéder en aucun cas aux données d'un autre client. Depuis cet espace, il peut créer son compte et vérifier son identité auprès des services de l'opérateur, afin de pouvoir ensuite agir avec l'agent.

\subsection{Supervision et administration}

L'utilisateur back-office doit disposer des moyens d'observer et de gouverner le système. Cela recouvre la consultation des conversations et de leurs transcriptions, l'inspection des décisions et des verdicts rendus, le suivi des indicateurs d'activité, et la gestion des conseillers et de leurs disponibilités. S'y ajoutent l'administration des règles métier et des référentiels, la gestion de la base documentaire, et le contrôle de l'intégrité du registre des opérations.

\subsection{Récapitulatif}

Les tableaux \ref{tab:bf} et \ref{tab:bf-suite} récapitulent ces besoins sous une forme identifiée, reprise au chapitre de réalisation pour établir la correspondance avec les mécanismes qui les mettent en oeuvre. Le premier couvre les domaines de la conversation, du contexte, de la connaissance et du début du traitement des actions. Le second poursuit sur les actions, les tickets, l'escalade, l'espace client, la supervision et l'administration.

\begin{table}[H]
\centering
\renewcommand{\arraystretch}{1.2}
\begin{tabular}{|p{1.3cm}|p{3.4cm}|p{9.2cm}|}
\hline
\multicolumn{3}{|c|}{\textbf{Besoins fonctionnels}} \\
\hline
\textbf{Réf.} & \textbf{Domaine} & \textbf{Énoncé} \\
\hline
BF01 & Conversation & Engager un échange parlé ou écrit et le mener à son terme sans contrainte de formulation. \\
\hline
BF02 & Conversation & Conserver le fil de l'échange d'un tour de parole au suivant. \\
\hline
BF03 & Conversation & Identifier l'intention exprimée par le client. \\
\hline
BF04 & Conversation & Conduire une clarification bornée à trois tentatives, en adaptant la stratégie, puis proposer un conseiller au delà. \\
\hline
BF05 & Conversation & Accepter une interruption du client à tout moment et reprendre sans perdre l'état. \\
\hline
BF06 & Conversation & Changer de langue en cours d'échange à la demande du client. \\
\hline
BF07 & Contexte & Reconstituer le profil, l'abonnement et la situation de facturation du client. \\
\hline
BF08 & Contexte & Récupérer l'historique des demandes et des échanges antérieurs. \\
\hline
BF09 & Contexte & Solliciter chaque élément de contexte uniquement lorsque la demande le justifie. \\
\hline
BF10 & Connaissance & Répondre aux questions documentaires à partir de la documentation de l'opérateur. \\
\hline
BF11 & Connaissance & Signaler explicitement l'absence de réponse plutôt que produire une réponse vraisemblable. \\
\hline
BF12 & Connaissance & Permettre l'alimentation et la mise à jour de la base documentaire par l'administrateur. \\
\hline
BF13 & Action & Vérifier l'identité du client avant toute opération sensible. \\
\hline
BF14 & Action & Produire un verdict explicite : autorisée, refusée, ou à transmettre à un conseiller. \\
\hline
\end{tabular}
\caption{Besoins fonctionnels de la plateforme}
\label{tab:bf}
\end{table}

\newpage

\begin{table}[H]
\centering
\renewcommand{\arraystretch}{1.2}
\begin{tabular}{|p{1.3cm}|p{3.4cm}|p{9.2cm}|}
\hline
\multicolumn{3}{|c|}{\textbf{Besoins fonctionnels (suite)}} \\
\hline
\textbf{Réf.} & \textbf{Domaine} & \textbf{Énoncé} \\
\hline
BF15 & Action & Associer à chaque verdict l'identifiant de la règle appliquée et sa justification. \\
\hline
BF16 & Action & Exécuter une action autorisée une fois et une seule. \\
\hline
BF17 & Action & Confirmer verbalement au client le résultat de l'opération. \\
\hline
BF18 & Ticket & Consulter les tickets existants d'un client. \\
\hline
BF19 & Ticket & Créer, mettre à jour et clôturer un ticket. \\
\hline
BF20 & Ticket & Tenir compte de l'historique des tickets pour éviter les doublons. \\
\hline
BF21 & Escalade & Reconnaître les situations appelant une intervention humaine. \\
\hline
BF22 & Escalade & Transférer la conversation à un conseiller disponible. \\
\hline
BF23 & Escalade & Négocier et réserver un créneau de rappel selon les disponibilités réelles. \\
\hline
BF24 & Escalade & Transmettre au conseiller le dossier complet de l'échange. \\
\hline
BF25 & Espace client & Consulter profil, facturation, demandes, rappels et historique des conversations. \\
\hline
BF26 & Espace client & Consulter et effacer les informations mémorisées par la plateforme. \\
\hline
BF27 & Espace client & Interdire l'accès aux données d'un autre client. \\
\hline
BF28 & Supervision & Consulter les conversations et leurs transcriptions. \\
\hline
BF29 & Supervision & Inspecter les décisions et les verdicts rendus. \\
\hline
BF30 & Supervision & Suivre les indicateurs d'activité de la plateforme. \\
\hline
BF31 & Supervision & Gérer les conseillers, leurs plannings et leurs indisponibilités. \\
\hline
BF32 & Administration & Administrer les règles métier et les référentiels. \\
\hline
BF33 & Administration & Gérer la base documentaire et ses versions. \\
\hline
BF34 & Administration & Contrôler l'intégrité du registre des opérations. \\
\hline
BF35 & Facturation & Ouvrir un litige sur une facture et le clôturer, la facture étant retirée du règlement pendant l'instruction. \\
\hline
BF36 & Espace client & Créer son compte depuis le portail et vérifier son identité auprès des services de l'opérateur, afin d'agir avec l'agent. \\
\hline
\end{tabular}
\caption{Besoins fonctionnels de la plateforme (suite)}
\label{tab:bf-suite}
\end{table}

\section{Besoins non fonctionnels}

Les besoins non fonctionnels expriment les propriétés de qualité attendues du système. Nous avons écarté les formulations générales pour ne retenir que les exigences réellement structurantes pour ce projet, chacune étant rattachée à la propriété qu'elle impose à l'architecture.

\subsection{Performance et comportement temps réel}

C'est l'exigence dominante. Un échange parlé n'est perçu comme naturel que si la réponse survient dans le rythme d'une conversation ordinaire. Cette exigence se décline en un budget de temps réparti entre les étapes du traitement : détection de fin de tour de parole, transcription, raisonnement, récupération documentaire, synthèse vocale et acheminement. Aucune étape ne peut être considérée isolément, puisque le dépassement de l'une consomme la marge de toutes les autres.

\subsection{Fiabilité et disponibilité}

La défaillance d'un fournisseur externe ne doit pas interrompre une conversation en cours. Le système doit prévoir, pour chaque étape dépendant d'un service tiers, une solution de repli capable de prendre le relais sans que le client en perçoive autre chose qu'une variation de qualité.

\subsection{Robustesse et tolérance aux pannes}

Le système doit continuer à fonctionner de manière dégradée mais correcte lorsqu'un composant non essentiel à l'échange devient indisponible. En particulier, une défaillance de la couche d'enregistrement ne doit jamais compromettre la conversation elle même. À l'inverse, aucune défaillance ne doit être silencieuse : elle doit rester visible pour le diagnostic.

\subsection{Sécurité et confidentialité}

L'accès aux fonctions du système doit être conditionné à une authentification, puis à une autorisation fondée sur le rôle. Un client ne doit pouvoir accéder qu'à ses propres données. Les opérations sensibles doivent exiger une vérification d'identité distincte de la simple ouverture de session. Les données personnelles doivent être masquées partout où leur valeur exacte n'est pas nécessaire. En l'absence d'information d'autorisation exploitable, le système doit refuser plutôt qu'autoriser.

\subsection{Traçabilité et auditabilité}

Tout acte ayant une conséquence pour un client doit laisser une trace exploitable indiquant qui a agi, sur quoi, à quel moment, avec quel résultat et en vertu de quelle règle. Ces traces doivent être organisées de telle sorte qu'une modification rétroactive soit détectable et non simplement découragée.

\subsection{Extensibilité et maintenabilité}

Le remplacement d'un fournisseur de transcription, de synthèse vocale ou de modèle de langue ne doit pas se propager au delà de la couche qui l'intègre. L'ajout d'un domaine métier ou d'une nouvelle capacité ne doit pas imposer la révision des composants existants. Les règles métier doivent rester indépendantes de toute bibliothèque tierce.

\subsection{Passage à l'échelle}

La plateforme doit prendre en charge plusieurs conversations simultanées, chacune disposant de son propre état, sans interférence entre elles. Les opérations concurrentes menées au sein d'une même conversation doivent être isolées les unes des autres.

\subsection{Interopérabilité et portabilité}

Les échanges avec les systèmes de l'opérateur et avec le système de gestion de tickets doivent reposer sur des interfaces standard. L'ensemble doit être déployable sous forme de conteneurs, indépendamment du système hôte.

\subsection{Testabilité}

Les composants portant une garantie du système doivent pouvoir être vérifiés isolément et de manière reproductible. Cette exigence conditionne directement la conception : elle impose que la logique de décision soit exempte d'effets de bord.

\subsection{Ergonomie}

L'interaction principale étant vocale, la qualité perçue tient à la nature de l'échange : formulations brèves, absence d'énumérations lues à voix haute, prise en compte immédiate des interruptions. Les interfaces destinées à l'utilisateur back-office doivent privilégier la lisibilité de l'information opérationnelle.

\subsection{Récapitulatif}

Le tableau \ref{tab:bnf} associe à chaque besoin non fonctionnel la propriété architecturale qui en découle.

\begin{table}[H]
\centering
\renewcommand{\arraystretch}{1.2}
\begin{tabular}{|p{1.3cm}|p{3.7cm}|p{8.9cm}|}
\hline
\textbf{Réf.} & \textbf{Exigence} & \textbf{Propriété imposée à l'architecture} \\
\hline
BNF01 & Performance temps réel & Budget de latence réparti par étape et mesuré de bout en bout. \\
\hline
BNF02 & Fiabilité et disponibilité & Chaîne de repli pour chaque dépendance externe de la chaîne vocale. \\
\hline
BNF03 & Robustesse & Découplage de l'enregistrement et de la conversation ; aucune défaillance silencieuse. \\
\hline
BNF04 & Sécurité et confidentialité & Authentification, autorisation par rôle, isolation des données, vérification d'identité renforcée, masquage, refus par défaut. \\
\hline
BNF05 & Traçabilité & Enregistrement de chaque acte conséquent et détection des altérations rétroactives. \\
\hline
BNF06 & Extensibilité & Dépendances externes confinées derrière des interfaces internes. \\
\hline
BNF07 & Maintenabilité & Règles métier indépendantes de toute bibliothèque tierce. \\
\hline
BNF08 & Passage à l'échelle & État par conversation et isolation des opérations concurrentes. \\
\hline
BNF09 & Interopérabilité & Interfaces standard vers les systèmes tiers. \\
\hline
BNF10 & Portabilité & Déploiement conteneurisé de l'ensemble des composants. \\
\hline
BNF11 & Testabilité & Logique de décision déterministe et exempte d'effets de bord. \\
\hline
BNF12 & Ergonomie & Réponses brèves adaptées à l'oral et prise en compte immédiate des interruptions. \\
\hline
\end{tabular}
\caption{Besoins non fonctionnels et propriétés architecturales associées}
\label{tab:bnf}
\end{table}

\section{Diagrammes de cas d'utilisation}

\subsection{Diagramme de cas d'utilisation général}

Le diagramme de la figure \ref{fig:cu-general} présente une vue d'ensemble des cas d'utilisation de la plateforme et de leur répartition entre les acteurs. Le client y conduit l'échange, qu'il s'agisse d'obtenir une information, d'effectuer une opération sécurisée, de gérer son profil, de signaler un problème ou de demander une prise en charge humaine. L'utilisateur back-office, qui regroupe les profils internes, y supervise les interactions, administre la base de connaissances et les règles métier et traite les demandes escaladées. Les systèmes tiers, facturation, gestion de tickets et notification, apparaissent en acteurs secondaires des cas qui les sollicitent ; les rôles précis de chaque utilisateur interne sont détaillés dans la matrice des droits du chapitre suivant.

\figureSVGPleine{Diagramme de cas d'utilisation général de la plateforme}{cu-general}{fig-cu-general}

Trois cas d'utilisation structurent l'ensemble du système et méritent un traitement détaillé : obtenir une information, effectuer une opération sécurisée et être mis en relation avec un conseiller humain. Le premier couvre le trajet complet d'une question documentaire, de la parole du client jusqu'à la réponse ancrée dans la documentation. Le deuxième couvre l'exécution d'une opération engageant le compte du client, qui concentre l'essentiel des garanties de sûreté de la plateforme. Le troisième couvre le passage à un conseiller humain, qui constitue la frontière entre l'autonomie du système et l'intervention humaine. Chacun fait l'objet d'un diagramme, d'une description textuelle et d'un diagramme de séquence système décrivant son scénario nominal.

\subsection{Cas d'utilisation : recherche et génération d'une réponse fondée sur les connaissances}

\subsubsection{Diagramme du cas}

Ce cas d'utilisation, présenté à la figure \ref{fig:cu-doc}, décrit la situation la plus fréquente : le client pose une question portant sur les offres, les procédures ou les conditions de service, et attend une réponse exacte. La réponse est établie à partir de la base de connaissances de l'opérateur. Le système interprète la demande, tient compte du contexte des échanges précédents afin de comprendre une question qui les prolonge, recherche les passages pertinents, les fusionne et évalue leur pertinence, puis génère une réponse fondée sur les connaissances retenues et accompagnée de leurs sources.

\figureTikz{Diagramme du cas d'utilisation Recherche et génération d'une réponse fondée sur les connaissances}{cu-doc}{fig-cu-doc}

\subsubsection{Description textuelle}

Le tableau \ref{tab:cu-doc} détaille le déroulement de ce cas.

\begin{table}[H]
\centering
\renewcommand{\arraystretch}{1.25}
\begin{tabular}{|p{3.0cm}|p{11.5cm}|}
\hline
\multicolumn{2}{|c|}{\textbf{Cas d'utilisation : recherche et génération d'une réponse fondée sur les connaissances}} \\
\hline
\textbf{Acteur principal} & Client \\
\hline
\textbf{Acteurs secondaires} & Aucun \\
\hline
\textbf{Objectif} & Obtenir une réponse exacte à une question portant sur les offres, les procédures ou les conditions de service, en la fondant sur la base de connaissances de l'opérateur et en tenant compte du contexte de l'échange. \\
\hline
\textbf{Préconditions} & Une conversation est établie entre le client et la plateforme. La base de connaissances est disponible et indexée. \\
\hline
\textbf{Postconditions} & Le client a reçu soit une réponse fondée sur les connaissances retrouvées, soit l'indication explicite que l'information n'est pas disponible. L'échange est enregistré. \\
\hline
\textbf{Scénario nominal} &
1. Le client formule sa question à l'oral. \newline
2. Le système transcrit la parole et interprète la demande, en tenant compte du contexte des échanges précédents. \newline
3. Le système recherche les passages pertinents dans la base de connaissances, puis dans la base vectorielle. \newline
4. Le système fusionne les passages candidats et évalue leur pertinence. \newline
5. Le système augmente la requête avec le contexte retenu et la transmet au moteur d'inférence. \newline
6. Le moteur d'inférence génère une réponse brève avec des références aux sources. \newline
7. Le système énonce la réponse au client et propose d'en obtenir le détail. \newline
8. L'échange est consigné dans le journal de conversation. \\
\hline
\textbf{Alternative A : demande ambiguë} & À l'étape 2, si la demande ne peut être interprétée avec certitude, le système pose une question de clarification unique, le client fournit la précision demandée, puis le traitement reprend à l'étape 3. \\
\hline
\textbf{Alternative B : absence d'information} & À l'étape 4, si aucun passage suffisamment pertinent n'est retrouvé, le système n'envoie aucune requête au moteur d'inférence : il indique au client qu'il ne dispose pas d'une information suffisante pour répondre et lui propose une prise en charge alternative. Le cas se termine. \\
\hline
\textbf{Alternative C : interruption} & À toute étape postérieure à 5, le client peut interrompre l'énoncé. Le système cesse de parler immédiatement et traite le nouveau tour de parole. \\
\hline
\end{tabular}
\caption{Description textuelle du cas Recherche et génération d'une réponse fondée sur les connaissances}
\label{tab:cu-doc}
\end{table}

\subsubsection{Diagramme de séquence système}

La figure \ref{fig:seq-doc} représente le scénario nominal sous la forme des échanges entre le client, l'agent spécialisé, la base de connaissances, la base vectorielle et le moteur d'inférence. Elle fait apparaître la prise en compte du contexte, la recherche en deux temps, la fusion et l'évaluation de la pertinence, puis la génération de la réponse avec ses références. Le fragment \texttt{opt} traduit la clarification en cas de demande ambiguë ; le fragment \texttt{alt} sépare le cas où la pertinence est confirmée de celui où aucun passage pertinent n'est retrouvé.

\figureTikz{Diagramme de séquence système du cas Recherche et génération d'une réponse fondée sur les connaissances}{seq-doc}{fig-seq-doc}

\subsection{Cas d'utilisation : réalisation d'une opération sécurisée}

\subsubsection{Diagramme du cas}

Ce cas, présenté à la figure \ref{fig:cu-action}, concentre les garanties de sûreté de la plateforme. Il couvre toute opération produisant un effet sur le compte du client, qu'il s'agisse d'un règlement, d'un report d'échéance, d'une recharge, d'un changement de forfait ou d'une opération sur la carte SIM ou sur la ligne.

La figure \ref{fig:cu-action} présente la réalisation d'une opération sécurisée. Avant toute exécution, le système vérifie l'identité du client par\,\acro{CIN}, valide l'opération au moyen des règles métier, la transmet au système métier concerné et l'enregistre dans le journal d'audit. Trois issues conditionnelles complètent ce trajet : restituer le résultat lorsqu'une opération a déjà été traitée, notifier le refus lorsque les conditions ne sont pas satisfaites, ou transmettre la demande à un conseiller lorsque l'opération ne peut être traitée automatiquement.

\figureTikz{Diagramme du cas d'utilisation Réalisation d'une opération sécurisée}{cu-action}{fig-cu-action}

\subsubsection{Description textuelle}

Le tableau \ref{tab:cu-action} détaille ce cas, qui est le plus contraint de la plateforme.

\begin{table}[H]
\centering
\renewcommand{\arraystretch}{1.25}
\begin{tabular}{|p{3.0cm}|p{11.5cm}|}
\hline
\multicolumn{2}{|c|}{\textbf{Cas d'utilisation : réalisation d'une opération sécurisée}} \\
\hline
\textbf{Acteur principal} & Client \\
\hline
\textbf{Acteurs secondaires} & Systèmes métier de l'opérateur \\
\hline
\textbf{Objectif} & Réaliser une opération sécurisée produisant un effet sur le compte du client, avec la garantie qu'elle est autorisée, exécutée une seule fois et consignée. \\
\hline
\textbf{Préconditions} & Une conversation est établie. Le client est rattaché à un compte. L'opération demandée figure parmi les opérations reconnues par le système. Une identité vérifiée par\,\acro{CIN} est exigée avant toute transmission au système métier. \\
\hline
\textbf{Postconditions} & L'opération a été exécutée et confirmée, ou refusée avec un motif, ou transmise à un conseiller. Dans tous les cas, un verdict motivé a été consigné dans le journal d'audit. \\
\hline
\textbf{Scénario nominal} &
1. Le client exprime la demande d'opération. \newline
2. Le système récapitule l'opération et sollicite une confirmation verbale explicite. \newline
3. Le client confirme. \newline
4. Le système vérifie l'identité du client par\,\acro{CIN}, si cette vérification n'est pas déjà valide. \newline
5. Le système constitue le contexte de la décision à partir de la situation vérifiée du client. \newline
6. Le système valide l'opération : il la soumet aux règles métier, qui rendent un verdict motivé. \newline
7. Le verdict étant favorable, le système transmet l'opération au système métier concerné, en garantissant qu'une reprise éventuelle ne produira pas de seconde exécution. \newline
8. Le système consigne l'opération et son verdict dans le journal d'audit. \newline
9. Le système confirme le résultat au client à l'oral. \\
\hline
\textbf{Alternative A : identité non confirmée} & À l'étape 4, si la vérification par\,\acro{CIN} échoue après le nombre d'essais autorisé, l'opération est abandonnée, le client en est informé, et le refus est consigné. \\
\hline
\textbf{Alternative B : opération refusée} & À l'étape 6, si les conditions de l'opération ne sont pas satisfaites, le système notifie le refus en énonçant au client le motif du refus dans des termes compréhensibles. L'opération n'est pas transmise. Le verdict est consigné. \\
\hline
\textbf{Alternative C : transmission à un conseiller} & À l'étape 6, si l'opération ne peut être traitée automatiquement, le système engage le cas d'utilisation d'escalade en transmettant le contexte constitué. \\
\hline
\textbf{Alternative D : opération déjà traitée} & Si la même opération est présentée une seconde fois pour la même conversation, le système restitue le résultat enregistré de la première exécution sans exécuter à nouveau. \\
\hline
\end{tabular}
\caption{Description textuelle du cas Réalisation d'une opération sécurisée}
\label{tab:cu-action}
\end{table}

\subsubsection{Diagramme de séquence système}

La figure \ref{fig:seq-action} représente la réalisation d'une opération sécurisée sous la forme des échanges entre le client, l'agent spécialisé, le moteur de politique, le système métier et le journal d'audit. Elle fait apparaître la confirmation verbale, la vérification d'identité par\,\acro{CIN} (fragment \texttt{opt}), puis la décision d'autorisation persistée avant toute exécution (fragment \texttt{alt} à trois branches : autorisée, refusée, traitement humain requis). Le fragment imbriqué d'idempotence garantit qu'une opération déjà traitée restitue son résultat sans seconde exécution. Aucune opération n'atteint le système métier sans un verdict favorable préalable, persisté et identifié.

\figureTikz{Diagramme de séquence système du cas Réalisation d'une opération sécurisée}{seq-action}{fig-seq-action}

\subsection{Cas d'utilisation : escalade vers un conseiller humain}

\subsubsection{Diagramme du cas}

Ce cas, présenté à la figure \ref{fig:cu-escalade}, décrit le passage de la prise en charge automatique à la prise en charge humaine. Il peut être déclenché par une règle métier, par une demande explicite du client ou par la constatation d'un mécontentement persistant au fil de la conversation. La figure \ref{fig:cu-escalade} représente l'escalade vers un conseiller humain : le système constitue le dossier de l'échange, recherche un conseiller disponible, puis transfère l'appel ou planifie un rappel. Dans ce dernier cas, il crée un ticket de rappel auprès du système de gestion de tickets et notifie le conseiller assigné par le service de notification. Le dossier complet, constitué dans les deux cas, est remis au conseiller avec le transfert.

\figureTikz{Diagramme du cas d'utilisation Escalade vers un conseiller humain}{cu-escalade}{fig-cu-escalade}

\subsubsection{Description textuelle}

Le tableau \ref{tab:cu-escalade} détaille ce cas.

\begin{table}[H]
\centering
\renewcommand{\arraystretch}{1.25}
\begin{tabular}{|p{3.0cm}|p{11.5cm}|}
\hline
\multicolumn{2}{|c|}{\textbf{Cas d'utilisation : escalade vers un conseiller humain}} \\
\hline
\textbf{Acteur principal} & Client \\
\hline
\textbf{Acteurs secondaires} & Conseiller, service de notification \\
\hline
\textbf{Objectif} & Confier la demande à un conseiller humain, immédiatement si possible, à une échéance convenue sinon, sans perte du contexte déjà constitué, et en conservant une trace de la demande auprès du système de suivi. \\
\hline
\textbf{Préconditions} & Une conversation est établie. Une condition d'escalade est satisfaite. \\
\hline
\textbf{Postconditions} & La conversation a été transférée à un conseiller, ou un rappel a été réservé sur un créneau réel ; dans ce cas un ticket de suivi a été créé auprès du système de gestion de tickets et le client a été notifié du créneau retenu. Le dossier de l'échange est disponible dans les deux cas. \\
\hline
\textbf{Scénario nominal} &
1. Le système constate la condition d'escalade et l'annonce au client. \newline
2. Le système constitue le dossier de l'échange : demande exprimée, éléments compris, opérations tentées et verdicts rendus. \newline
3. Le système recherche un conseiller disponible correspondant au domaine de la demande. \newline
4. Un conseiller étant disponible, le système le réserve. \newline
5. Le système transfère l'appel vers ce conseiller. \newline
6. Le conseiller prend en charge la conversation, muni du dossier. \\
\hline
\textbf{Alternative A : aucun conseiller disponible} & À l'étape 3, si aucun conseiller n'est disponible, le système en informe le client et lui propose un rappel. Il présente des créneaux issus des disponibilités réelles, convient d'un créneau avec le client, le réserve, crée un ticket de suivi et notifie le client du créneau retenu. Le conseiller concerné est avisé avec le dossier. \\
\hline
\textbf{Alternative B : échec du transfert} & À l'étape 5, si le transfert échoue, le conseiller réservé est libéré, le client en est informé sans ambiguïté, et le système bascule sur la proposition de rappel décrite en alternative A. \\
\hline
\textbf{Alternative C : refus du rappel} & Si le client décline la proposition de rappel, le système clôt l'échange en récapitulant ce qui a été fait et en indiquant les moyens de contact disponibles, sans créer de rappel. \\
\hline
\end{tabular}
\caption{Description textuelle du cas Escalade vers un conseiller humain}
\label{tab:cu-escalade}
\end{table}

\subsubsection{Diagramme de séquence système}

La figure \ref{fig:seq-escalade} représente le scénario nominal et son alternative principale sous la forme des échanges entre le client, l'agent de supervision, le service de routage, le système de tickets, le service de notification et le conseiller. L'agent de supervision constitue le dossier et recherche un conseiller disponible. Le fragment \texttt{alt} sépare le cas où un conseiller est disponible, avec transfert de l'appel et remise du dossier, de celui où aucun conseiller n'est disponible ou où le transfert échoue, qui conduit à la création d'un ticket de rappel, à la notification du conseiller assigné et à la confirmation orale de l'horaire retenu.

\figureTikz{Diagramme de séquence système du cas Escalade vers un conseiller humain}{seq-escalade}{fig-seq-escalade}

\section*{Conclusion}

Ce chapitre a transformé les objectifs du projet en spécifications exploitables. Nous avons identifié deux acteurs principaux, le client et l'utilisateur back-office, ainsi que trois familles d'acteurs secondaires correspondant aux systèmes que la plateforme sollicite sans les posséder. Les responsabilités propres à l'administrateur, au superviseur et au conseiller sont détaillées dans la matrice des droits du chapitre suivant. Nous avons ensuite formulé trente quatre besoins fonctionnels répartis en huit domaines, puis douze besoins non fonctionnels dont chacun impose une propriété identifiée à l'architecture. Les cas d'utilisation ont enfin été présentés dans leur ensemble, et détaillés pour les trois qui structurent le système : la recherche et la génération d'une réponse fondée sur les connaissances, la réalisation d'une opération sécurisée et l'escalade vers un conseiller humain.

Ces trois cas font apparaître une même exigence sous des formes différentes. La recherche et la génération d'une réponse fondée sur les connaissances imposent que le système sache reconnaître les limites de ce qu'il sait. La réalisation d'une opération sécurisée impose qu'il sache reconnaître les limites de ce qu'il a le droit de faire. L'escalade vers un conseiller humain impose qu'il sache reconnaître les limites de ce qu'il peut traiter seul. Le chapitre suivant montre comment l'architecture de la solution répond à ces trois exigences, en présentant d'abord la vue statique du système, puis sa vue dynamique, avant de détailler la conception de la couche agentique et celle du contrôle d'accès.
